\section{最终识别}\label{sec:final-identifications}

MWW 定理~3.10 将迭代商与三柄后模等同，命题~3.4 给出四柄同构 \cite{ManolescuWalkerWedrich2023}。命题~\ref{thm:P0discharge}、\ref{thm:Cdischarge} 与~\ref{thm:Sdischarge} 记录假设~\ref{hyp:P0}--\ref{hyp:P2} 的状态；三者均为 \Open。对 Johnson 替换四柄图景 $X_J$，三个 $3$-柄沿反向 $1$-柄的 belt 球面附着后恢复边界 $S^3$（假设~\ref{hyp:P3}），附着 PL $4$-ball，标准球面模的 quantum degree-$494$ 分项由假设~\ref{hyp:P3} 状态注记中的计算连同 MWW 推论~3.5 而消失。一旦假设~\ref{hyp:P0} 确立，分阶段 Kirby 管线将把 $X_J$ 与 $\Sigma_A^0$ 识别。

\begin{remark}[形式化边界]\label{rem:lean-boundary}
Lean 开发 \cite{Lean4} 将定理~\ref{thm:A} 证明为从结构 \texttt{ExternalGeometry} 与 \texttt{CSExternalGeometry} 出发的条件蕴含。它检验矩阵算术、值 $2624$、次数 $494$ 与抽象商引理，但不构造本文证明的几何实例。附录~\ref{app:lean-fields} 列出与 Lean 字段的对应。
\end{remark}
